%%%%%%%%%%%%%%%%%%%%%%%%%%%%%%%%%%%%%%%%%%%%%%%%%%%%%%%%%%%%
\section{Introduction}
\label{sec:intro}
%%%%%%%%%%%%%%%%%%%%%%%%%%%%%%%%%%%%%%%%%%%%%%%%%%%%%%%%%%%%

Foundation models pretrained on large unlabeled scientific corpora are reshaping prediction across the natural sciences. Protein language models such as ESM-2~\citep{lin2023esm2} learn evolutionary patterns that support fold and variant-effect prediction, and AlphaFold~\citep{jumper2021highly} predicts atomic protein structure at near-experimental accuracy. Genomic foundation models such as Evo~\citep{nguyen2024evo} model regulatory and functional DNA at base-pair resolution. Global weather models such as GraphCast~\citep{lam2023graphcast} now match or surpass numerical weather prediction on ten-day forecasts. Across these domains, large transformers pretrained on raw scientific data help unlocked downstream tasks that previously required hand-engineered features or domain-specific inductive biases.

Similarly, chemisrty models pretrained on hundreds of millions to billions of SMILES strings, including ChemBERTa~\citep{chithrananda2020chemberta}, \molm~\citep{ross2022largescalechemicallanguagerepresentations}, \smi~\citep{soares2024smi}, and ChemFM~\citep{cai2026chemfm}, now serve as default backbones for molecular property prediction across MoleculeNet~\citep{wu2018moleculenet}, ADMET screening~\citep{harod2025effichem}, and multi-scale materials and battery design~\citep{sharma2025foundation}. Like their counterparts in protein and genomic modeling, these encoders learn molecular representations from raw text without explicit structural inductive biases, and yet underwrite strong performance on a broad range of downstream chemistry tasks.

Strong downstream performance, however, does not imply that we understand what these models have learned. Because they are typically deployed as black-box feature extractors, basic questions about their internal representational mechanisms remain open: do attention heads track specific chemical substructures? at which depth does a contextual property such as ring membership or chirality become linearly recoverable? are those representations actually used by the encoder, or only co-encoded with the eventual prediction? Parallel work in other scientific domains shows why such questions matter beyond curiosity. \citet{lu2026mechanisms} use counterfactual interventions on ESMFold to separate spatial-structure latents from purely biochemical ones, exposing concrete intervention points for diagnosing folding behavior. \citet{goodfire2026evee} convert Evo~2 embeddings into state-of-the-art interpretable variant-effect predictions for the 4.2 million variants in ClinVar. \citet{goodfire2026alzheimers} apply sparse autoencoders to an epigenetic foundation model and find that DNA fragment-length patterns — not methylation, the assumed signal — drive Alzheimer's classification, yielding a new biomarker class identified by interpretation rather than \emph{de novo} discovery. Going further back in the ML literature, \citet{mcgrath2022acquisition} probe AlphaZero and recover human chess concepts that the model independently rediscovered through self-play. Taken together, these results articulate two concrete payoffs of interpreting scientific foundation models: \emph{(i)} when a learned representation is wrong or shortcuts onto a confounder, identifying it suggests how to make the next model better; \emph{(ii)} when a learned representation is right, decoding it can surface scientific structure that was not explicitly built in.

Motivated by these payoffs, we carry the same program over to two of the most-used SMILES encoders, \molm{} and \smi{}. The two models are similar at the architectural level — same depth (12), hidden size (768), and linear-attention backbone — so representational differences can plausibly be attributed to pretraining design rather than architecture, which makes them a useful pair for asking what role pretraining plays in the resulting representations. We combine three complementary analyses on each: (1) a layerwise attention--distance correlation across all heads, asking whether attention tracks 3D molecular geometry; (2) layerwise linear probing for eight per-atom RDKit-derived chemical properties, separating what is recoverable from SMILES tokenization from what requires contextual computation; (3) causal interventions along the probe directions, which test whether a decoded property is functionally used by the rest of the encoder rather than merely linearly accessible.

Our main findings are:
\begin{itemize}
    \item \textbf{Attention only weakly tracks 3D geometry in either model.} Mean Spearman correlation between attention and inverse interatomic distance stays below $\rho \approx 0.3$ across all layers and heads of both models, with \smi{} marginally stronger than \molm{}. This contrasts with prior \molm{} claims that pooled across heads before computing alignment.
    \item \textbf{Surface-form chemistry is encoded at the embedding; contextual chemistry is not.} Properties exposed by SMILES tokenization (atom type, aromaticity) reach $100\%$ probe accuracy at layer~0 in both models. Properties that require resolving context (ring membership, degree, chiral tag, hybridization, valence) only become linearly recoverable inside the encoder.
    \item \textbf{\smi{} encodes contextual properties more cleanly than \molm{}.} \smi{} reaches near-perfect probe accuracy by mid-depth on ring membership, degree, and chiral tag, while \molm{} plateaus well below --- and sits near the majority-class baseline on chiral tag. The gap persists under a nonlinear probe.
    \item \textbf{The two models use the same properties at different depths.} Causal-intervention success along probe directions reaches $\approx 1.0$ from mid-layers in \smi{}, but only near the top of the encoder in \molm{}, even where \molm{}'s probe accuracy is already high. Linear recoverability and causal use are dissociable, and the depth profile is itself a signature of pretraining design.
\end{itemize}
